% ch09-risks.tex – Chapter 9: risk analysis and mitigation plan (final round, Writer). Register rows generated from data/risks.csv.
\ChapterMission{המשימה: לזהות מה עלול להשתבש, לדרג את הסיכונים ולהכין לכל סיכון קריטי הפחתה ותוכנית חלופית}
\ChapterIcon{warning}
\chapter{ניתוח סיכונים ותוכנית הפחתה}\label{chap:risks}
\begin{scenebox}
הרמזור מתקצר, והאיש הירוק נדלק מוקדם מן הרגיל. ילד שעמד על המדרכה כבר יודע את זה:
בפעם הבאה הוא עומד עם רגל אחת בכביש, מוכן לזנק עוד לפני שהאור מתחלף.
\end{scenebox}

הילד שלמד לצפות את השינוי ממחיש שגם מערכת שנועדה לשפר בטיחות יוצרת סיכונים חדשים.
הפרק מזהה את הסיכונים, מדרג אותם ומגדיר לכל סיכון קריטי צעדי הפחתה ותוכנית חלופית.

\section{שיטת הדירוג}
לכל סיכון נקבעו הסבירות (\en{P}) וההשפעה (\en{I}) בסולם \range{1}{5}, ודרגת הסיכון היא מכפלתן (משוואה~\eqnum{eq:risk}):
\begin{equation}
  R = P \times I, \qquad R \ge 12 \;\Rightarrow\; \hebmath{סיכון קריטי}
  \label{eq:risk}
\end{equation}
הערכים הם הערכות הצוות, ולכל סיכון נקבעה גם דרגה שיורית, אחרי ההפחתה.

\section{מרשם הסיכונים}
טבלה~\ref{tab:risks} מציגה את שמונת הסיכונים. ארבעה מהם קריטיים, וכולם נובעים מן החידוש עצמו: האישור לקיצור האדום, דיוק החישה, הפרטיות והתגובה של הולכי הרגל לשינוי.
איור~\ref{fig:g4} ממקם את הסיכונים במפת סבירות והשפעה, ומראה שאף סיכון אינו נשאר באזור הקריטי אחרי ההפחתה.

\begin{table}[ht]
  \centering
  \caption{מרשם הסיכונים: סבירות, השפעה, דרגה לפני ההפחתה ודרגה שיורית (הערכות הצוות; מודגש = קריטי)}
  \label{tab:risks}
  \small
  \begin{tabularx}{\linewidth}{lXlcccc}
    \toprule
    \headrow \thead{מזהה} & \thead{תרחיש} & \thead{תחום} & \thead{\en{P}} & \thead{\en{I}} & \thead{\en{R}} & \thead{שיורי} \\
    \midrule
    \en{R1} & אישור האדום המסתגל מתעכב או נדחה & רגולציה & \num{3} & \num{5} & \textbf{\num{15}} & \num{8} \\
    \en{R2} & החמצה או זיהוי שווא של ממתינים & טכנולוגיה & \num{4} & \num{3} & \textbf{\num{12}} & \num{6} \\
    \en{R3} & פגיעה בפרטיות בשל מצלמה & פרטיות & \num{3} & \num{4} & \textbf{\num{12}} & \num{4} \\
    \en{R4} & עיכוב נוסף לכלי רכב & תנועה & \num{3} & \num{3} & \num{9} & \num{6} \\
    \en{R5} & יציאות מוקדמות של הולכי רגל & התנהגות & \num{3} & \num{4} & \textbf{\num{12}} & \num{8} \\
    \en{R6} & קמפיין ללא השפעה & התנהגות & \num{3} & \num{2} & \num{6} & \num{4} \\
    \en{R7} & תזמון הספירה אינו מתואם עם הבקר & טכנולוגיה & \num{2} & \num{4} & \num{8} & \num{4} \\
    \en{R8} & עיכוב ברכש ובהתקנה & ניהול & \num{3} & \num{3} & \num{9} & \num{6} \\
    \bottomrule
  \end{tabularx}
\end{table}

\begin{figure}[ht]
  \centering
  \includefigure{figures/g4-risk-heatmap.pdf}
  \caption{מפת סבירות והשפעה של הסיכונים; האזור הקריטי מסומן בצבע ובדוגמה (הערכות הצוות)}
  \label{fig:g4}
\end{figure}

\section{הפחתה ותוכנית חלופית לסיכונים הקריטיים}
\begin{description}
\item[אישור רגולטורי (\en{R1}).] \textbf{הפחתה:} תיאום עם העירייה כבר בחבילת הייזום, ותכנון שאינו נוגע ברצפות התקן \cite{fhwa2025mutcd}. \textbf{חלופה:} ספירה בפינוי וקמפיין בלבד, כלומר \en{A3}, עד לקבלת ההיתר.
\item[החמצה או זיהוי שווא (\en{R2}).] דיוק החישה באב-טיפוס היה בינוני \cite{wickramasinghe2020}, וגלאים דרשו כיול \cite{hughes2000}. \textbf{הפחתה:} בדיקת מעבדה, כיול ואימות מול וידאו מתויג (\en{REQ-02}). \textbf{חלופה:} חזרה אוטומטית לתזמון הקבוע (\en{REQ-10}).
\item[פרטיות (\en{R3}).] \textbf{הפחתה:} עיבוד בקצה, השבתת הזרמת הווידאו \cite{agd2024agd645} וספירות בלבד (\en{REQ-11}). \textbf{חלופה:} גלאי שאינו מצלם, מכ״ם או אינפרה-אדום.
\item[יציאות מוקדמות (\en{R5}).] ספירה בזמן האדום כבר נקשרה ליותר יציאות מוקדמות של ילדים \cite{fu2016}. \textbf{הפחתה:} אין ספירה באדום, ומועד הקיצור אינו מוצג מראש; היציאות המוקדמות נמדדות בפיילוט. \textbf{חלופה:} השבתת הקיצור אם שיעורן עולה.
\end{description}
לשאר הסיכונים די בהפחתה: תקרה לקיצור ושמירה על הירוק המזערי לעיכוב הרכב (\en{R4}), בדיקה מוקדמת של מסרי הקמפיין (\en{R6}), תזמון הספירה מן הבקר (\en{R7}), והזמנה מוקדמת של הציוד (\en{R8}).
